% =====================================================================
% appendix.tex — STUB.
%
% Filled in across multiple manuscript increments:
%   RB-003 / sec:appendix-deriv-a1     — A1 ρ-channel derivation
%   RB-026 / sec:appendix-deriv-c2     — independent r†(v) re-derivation
%                                        (+ ρ>0 extension when ready)
%   RB-013 / sec:appendix-c5           — symmetric-recovery at r=1 note
%
% The §3.3 confident-spine results (C2 non-monotonicity, C5 symmetric
% recovery, central-tendency C1) get their full derivations here;
% extension levers (A1 ρ, A3 conservation family, A2/A8 heterogeneity)
% each get their own appendix subsection.
% =====================================================================

\appendix
\section{Derivations and consistency checks}
\label{sec:appendix}

% ---------------------------------------------------------------------
\subsection{A1 decorrelation-channel derivation}
\label{sec:appendix-deriv-a1}

[Filled by RB-003 (derivation increment). Independent re-derivation
of the equicorrelated-Gaussian orthant integral
$\PnoFA(\corr) = \int \phi(z)\prod_i \Phinorm\!\big((c_i +
\sqrt{\corr}\,z)/\sqrt{1 - \corr}\big)\, dz$,
one-factor Gauss--Hermite reduction, and the Slepian-monotonicity
argument that bounds $\partial \CF / \partial \corr$ in the
variant-A case. Cites~\cite{Slepian1962}.]

% ---------------------------------------------------------------------
\subsection{Closed-form VDA escape threshold $\rdagger(\val;\corr)$}
\label{sec:appendix-deriv-c2}

This subsection promotes the closed-form escape threshold introduced
in Section~\ref{sec:results-c2} into a derivation in the rebuild's
voice and extends it from the inherited $\corr = 0$ case to the
$\corr \in [0, 1)$ family the A1 lever
(Section~\ref{sec:model-rho-channel}) opens. Every numerical
statement here traces to the companion verification script
\texttt{Rebuild/derivations/verify\_C2\_rho/verify.py}, whose
output digest is recorded at the end of the subsection.

\paragraph{Setup at the boundary.}
At the uniform allocation $\alpha = 1/\Nloc$ the asymmetric branch
of the sensitivity rule of Section~\ref{sec:model} collapses to
$\dprime_c = \dprime_u = \dprime_{\mathrm{base}} :=
\dprimemax\,f(1/\Nloc)$ irrespective of $(\Rsens, \corr)$: the
slope brackets vanish at $\alpha = 1/\Nloc$ before being weighted by
$\benefit(\Rsens)$ or $\cost(\Rsens)$. The criterion optimum of
policy P3 at $\alpha = 1/\Nloc$ is therefore the joint argmax of the
expected reward over a fixed-$\dprime$ landscape that depends on
$\corr$ only through the correlated no-FA term $\PnoFA(\corr)$:
\begin{equation}
  \bigl(c_c^\star(\corr),\,c_u^\star(\corr)\bigr)
  \;:=\;
  \operatorname*{arg\,max}_{(c_c, c_u)}\,
  \Rew\bigl(\dprime_{\mathrm{base}},\,\dprime_{\mathrm{base}},\,
           c_c,\,c_u;\,\valid,\val,\Nloc,\CR,\corr\bigr).
  \label{eq:c2-boundary-crit}
\end{equation}
At the C2 headline cell
$(\Nloc, \valid, \val) = (4, 0.5, 5)$, $\CR = 3.0$ (variant~A) the
$\corr = 0$ optimum $(c_c^\star, c_u^\star) = (0.10, 1.75)$ shifts
to $(0.05, 1.80)$ at $\corr = 0.2$: the cued criterion drops by
$0.05$ (more liberal, because a cued FA is partially
``explained away'' by correlated FAs at the other locations) and
the uncued criterion rises by $0.05$ (more conservative, because
an uncued FA now also hurts the cued no-FA joint event).

\paragraph{The $\corr$-aware $d$-gradient integrals.}
Differentiation of the correlated orthant probability
$\PnoFA(\corr)$ (Equation~\eqref{eq:pnofa-rho}) under the integral
sign --- both the integrand and its $\dprime$-derivatives are
bounded, so dominated convergence applies --- gives the two
$\corr$-aware 1-D Gauss--Hermite integrals
\begin{align}
  I_c(b_c, b_u; \corr, \Nloc)
  &\;:=\;
  \int_{-\infty}^{\infty}
  \frac{1}{\sqrt{1 - \corr}}\,
  \phi\!\left(\frac{b_c - \sqrt{\corr}\,z}{\sqrt{1 - \corr}}\right)
  \Phinorm\!\left(\frac{b_u - \sqrt{\corr}\,z}{\sqrt{1 - \corr}}\right)^{\Nloc - 1}
  \phi(z)\,dz,
  \label{eq:c2-Ic}
  \\[2pt]
  I_u(b_c, b_u; \corr, \Nloc)
  &\;:=\;
  \int_{-\infty}^{\infty}
  \frac{1}{\sqrt{1 - \corr}}\,
  \Phinorm\!\left(\frac{b_c - \sqrt{\corr}\,z}{\sqrt{1 - \corr}}\right)
  \Phinorm\!\left(\frac{b_u - \sqrt{\corr}\,z}{\sqrt{1 - \corr}}\right)^{\Nloc - 2}
  \phi\!\left(\frac{b_u - \sqrt{\corr}\,z}{\sqrt{1 - \corr}}\right)
  \phi(z)\,dz,
  \label{eq:c2-Iu}
\end{align}
where $b_i := c_i + \dprime_{\mathrm{base}}/2$. Both reduce, at
$\corr = 0$, to the independent products
$I_c^0 = \phi(b_c)\,\Phinorm(b_u)^{\Nloc - 1}$ and
$I_u^0 = \Phinorm(b_c)\,\Phinorm(b_u)^{\Nloc - 2}\,\phi(b_u)$ that
appear in the $\corr = 0$ chain rule. In the rebuilt model
implementation, \eqref{eq:c2-Ic}--\eqref{eq:c2-Iu} share the same
64-node $(z_k, w_k)$ Gauss--Hermite nodes that
\texttt{Rebuild/model/core.py:p\_no\_fa\_grid} uses for
$\PnoFA(\corr)$ itself.

\paragraph{Boundary first-order condition.}
The chain-rule expansion of $\partial \Rew(\alpha)/\partial\alpha$
at $\alpha = 1/\Nloc^{+}$ collects the asymmetric P3 booking into
\begin{equation}
  \frac{\partial \Rew}{\partial\alpha}\bigg|_{1/\Nloc^{+}}\!(\corr)
  \;=\;
  \dprimemax\,f'(1/\Nloc)\,
  \Bigl[\,\benefit(\Rsens)\,K_c(\val; \corr)
        \;-\; \cost(\Rsens)\,\frac{K_u(\val; \corr)}{\Nloc - 1}\,\Bigr],
  \label{eq:c2-rho-foc}
\end{equation}
with $\corr$-aware coefficients
\begin{align}
  K_c(\val; \corr)
  &= \tfrac{1}{4}\Big[\valid\,\val\,
                      \phi\!\big(\tfrac{\dprime_b}{2} - c_c^\star(\corr)\big)
                      \;+\; \CR(\val)\,
                      I_c\bigl(b_c^\star(\corr), b_u^\star(\corr); \corr, \Nloc\bigr)\Big],
  \label{eq:c2-Kc-rho}
  \\[2pt]
  K_u(\val; \corr)
  &= \tfrac{1}{4}\Big[(1 - \valid)\,
                      \phi\!\big(\tfrac{\dprime_b}{2} - c_u^\star(\corr)\big)
                      \;+\; (\Nloc - 1)\,\CR(\val)\,
                      I_u\bigl(b_c^\star(\corr), b_u^\star(\corr); \corr, \Nloc\bigr)\Big],
  \label{eq:c2-Ku-rho}
\end{align}
where $b_i^\star(\corr) := c_i^\star(\corr) + \dprime_{\mathrm{base}}/2$
and $\dprime_b = \dprime_{\mathrm{base}}$ as before. The
change-trial density terms (the $\phi$ summands in
\eqref{eq:c2-Kc-rho}--\eqref{eq:c2-Ku-rho}) inherit $\corr$
dependence only through $c_i^\star(\corr)$; the no-FA terms (the
$I_c, I_u$ summands) carry the full $\corr$ dependence of the
correlated quadrature.

\begin{proposition}[Escape threshold under equicorrelation]
\label{prop:r-dagger-rho}
Fix $(\Nloc, \valid, \val, \dprimemax, f_0, h, \CR, \corr)$ with
$\corr \in [0, 1)$ and let $(c_c^\star(\corr), c_u^\star(\corr))$
solve \eqref{eq:c2-boundary-crit}. Then
\begin{equation}
  \rdagger(\val; \corr)
  \;=\;
  \frac{K_u(\val; \corr)}{(\Nloc - 1)\,K_c(\val; \corr)},
  \label{eq:r-dagger-rho}
\end{equation}
with $K_c(\val; \corr)$ and $K_u(\val; \corr)$ defined by
\eqref{eq:c2-Kc-rho}--\eqref{eq:c2-Ku-rho}. Below $\rdagger(\val;\corr)$,
$\alpha^\star_{\mathrm{P_1}}(\Rsens, \val; \corr) = 1/\Nloc$; above,
$\alpha^\star_{\mathrm{P_1}}(\Rsens, \val; \corr) > 1/\Nloc$.
\end{proposition}

\begin{proof}
Setting \eqref{eq:c2-rho-foc} to zero and dividing out the strictly
positive prefactor
$\dprimemax\,f'(1/\Nloc)\,\cost(\Rsens) > 0$ gives
$\Rsens\,K_c(\val;\corr) = K_u(\val;\corr)/(\Nloc - 1)$ (using
$\benefit(\Rsens)/\cost(\Rsens) = \Rsens$ from
$\benefit/\cost = \Rsens$, which holds in every member of the
conservation family of Section~\ref{sec:extensions-a3});
solving for $\Rsens$ gives \eqref{eq:r-dagger-rho}. A one-sided
$\alpha$ finite-difference at
$\Rsens = \rdagger(\val;\corr) \pm 0.05$ with
$\Delta\alpha = 10^{-3}$ confirms the sign of
$\partial \Rew/\partial\alpha|_{1/\Nloc^{+}}$ flips at every
$(\val, \corr) \in \{1, 2, 3\} \times \{0, 0.2\}$ probe
(\texttt{boundary\_FD\_check}, $6/6$ flips).
\end{proof}

\paragraph{Recovery at $\corr = 0$.}
At $\corr = 0$ the integrand of \eqref{eq:c2-Ic} loses its
$z$-dependence ($\sqrt{\corr} = 0$, $\sqrt{1 - \corr} = 1$,
$\int \phi(z)\,dz = 1$), so
$I_c(b_c, b_u; 0, \Nloc) = \phi(b_c)\,\Phinorm(b_u)^{\Nloc - 1} =
I_c^0$ and likewise
$I_u(b_c, b_u; 0, \Nloc) = \Phinorm(b_c)\,\Phinorm(b_u)^{\Nloc - 2}\,\phi(b_u) = I_u^0$.
The criterion optimum \eqref{eq:c2-boundary-crit} at $\corr = 0$
recovers the asymmetric P3 optimum of Section~\ref{sec:results-c2},
so $K_c(\val; 0) = K_c(\val)$ and $K_u(\val; 0) = K_u(\val)$ of
Equations~\eqref{eq:K-c}--\eqref{eq:K-u}, and
Proposition~\ref{prop:r-dagger-rho} reduces to
Equation~\eqref{eq:r-dagger}. The recovery is \emph{structural} ---
the integrals collapse term-by-term, not merely to numerical
tolerance. Numerically, the verification script computes
$|\rdagger_{\mathrm{rb\text{-}026}}(\val; 0) -
\rdagger_{\mathrm{rb\text{-}006}}(\val)| \le 5 \times 10^{-4}$
across $\val \in \{1, 2, 3, 5, 8, 10\}$; the residual is the
$\Delta c = 0.05$ quantisation of $(c_c^\star, c_u^\star)$ on the
shared criterion grid (at $\val = 2$ both implementations land on
the same grid optimum and \eqref{eq:r-dagger-rho} matches
Equation~\eqref{eq:r-dagger} to binary identity).

\paragraph{Drift prediction at the headline cell.}
The empirical drift table reported in Section~\ref{sec:results-c2}
as ``$r^\star$ drifts upward with $\corr$ at low $\val$''
(Table~\ref{tab:rho-sensitivity}) now has an analytic
counterpart: the \emph{lower edge} of the escape band
$\rdagger(\val;\corr)$, evaluated from
\eqref{eq:r-dagger-rho}, drifts upward at every $\val$ in the
family. Table~\ref{tab:r-dagger-rho-drift} reports the prediction.
%
\begin{table}[t]
  \centering
  \small
  \begin{tabular}{r r r r r c}
    \toprule
    $\val$ &
    $\rdagger(\val;0)$ &
    $\rdagger(\val;0.2)$ &
    $\Delta\rdagger$ &
    $\%\,\Delta\rdagger$ &
    sign-match $r^\star$? \\
    \midrule
    $1$  & $0.343$ & $0.354$ & $+0.0103$ & $+3.0\%$  & --- (P2 ref) \\
    $2$  & $0.168$ & $0.173$ & $+0.0051$ & $+3.0\%$  & \checkmark \\
    $3$  & $0.099$ & $0.113$ & $+0.0132$ & $+13.2\%$ & \checkmark \\
    $5$  & $0.050$ & $0.058$ & $+0.0077$ & $+15.2\%$ & \checkmark \\
    $8$  & $0.022$ & $0.029$ & $+0.0066$ & $+29.7\%$ & \checkmark \\
    $10$ & $0.016$ & $0.019$ & $+0.0030$ & $+18.7\%$ & \checkmark \\
    \bottomrule
  \end{tabular}
  \caption{Closed-form drift of the escape threshold under
  correlation at the C2 headline cell
  $(\Nloc, \valid, \dprimemax, f_0, h) = (4,\,0.5,\,2,\,0.5,\,\sqrt{})$,
  variant~A. $\rdagger(\val;\corr)$ from
  Equation~\eqref{eq:r-dagger-rho}. The drift
  $\Delta\rdagger := \rdagger(\val;0.2) - \rdagger(\val;0)$ is
  \emph{positive} at every $\val$ tested, including $\val = 1$
  (the value-blind P2 reference): the regime in which P1 is stuck
  at uniform allocation \emph{widens} as $\corr$ grows. The
  closed-form sign of $\Delta\rdagger$ matches the empirical peak
  drift sign $\Delta r^\star$ of
  Table~\ref{tab:rho-sensitivity} at all five $\val \ne 1$
  rows.}
  \label{tab:r-dagger-rho-drift}
\end{table}
%
Two findings.
%
\begin{itemize}
  \item \emph{Sign.} $\Delta\rdagger(\val) > 0$ at every $\val$
  tested. The escape-band lower edge moves up under correlation
  across the full $\val$-family; the lower edge cannot
  exceed the peak, so the peak must also drift up
  (Section~\ref{sec:results-c2}, ``$r^\star$ drifts up under
  $\corr$''). The five $\val \ne 1$ rows of
  Table~\ref{tab:r-dagger-rho-drift} match the sign of the empirical
  peak drift $\Delta r^\star$ reported in
  Table~\ref{tab:rho-sensitivity} ($5/5$ sign-match).
  \item \emph{Magnitude.} The percentage drift $\%\,\Delta\rdagger$
  is largest at $\val = 8$ ($+30\%$) and smallest at $\val = 1$
  ($+3\%$). The absolute drift $\Delta\rdagger$ is bounded by
  $\le 0.014$ across the family, while the empirical \emph{peak}
  drift $\Delta r^\star$ in Table~\ref{tab:rho-sensitivity}
  ranges up to $0.13$ at $\val = 2$. The peak drift exceeds the
  threshold drift because the peak also responds to the
  corresponding upward drift of $\rdagger(\val = 1; \corr)$ (which
  widens the escape band's \emph{upper} edge) and to changes inside
  the band as $\alpha^\star_{\mathrm{P_1}}$ climbs differently
  across $\corr$. Equation~\eqref{eq:r-dagger-rho} pins the
  lower-edge drift only; a closed form for the peak location
  $r^\star(\val;\corr)$ that would additionally model the
  $\alpha^\star$-climb above $\rdagger$ is left to a follow-up
  derivation increment.
\end{itemize}

\paragraph{Mechanistic mirror to A1.}
The two $\corr$-aware integrals $I_c, I_u$ of
\eqref{eq:c2-Ic}--\eqref{eq:c2-Iu} are the gradient analogue of
the orthant decomposition behind the A1 ``three levers, not two''
reframe of Section~\ref{sec:model-three-levers}. The same
one-factor Gauss--Hermite reduction that makes $\PnoFA(\corr)$
tractable in the model code makes the boundary slope
$\partial \Rew/\partial\alpha|_{1/\Nloc^{+}}$ tractable here.
Section~\ref{sec:results-c2} reports the peak drift
\emph{empirically} as a consequence of the A1 lever;
Proposition~\ref{prop:r-dagger-rho} makes the
\emph{lower-edge} part of that drift an analytic prediction of
the rebuilt model. What
Section~\ref{sec:appendix-deriv-a1}'s Slepian monotonicity argument
does for $\CF$ as a function of $\corr$,
Proposition~\ref{prop:r-dagger-rho} does for the
$\rdagger$ component of the escape band: the rebuild now has a
closed-form locus for one of the two channels of
$\partial\VDA/\partial\corr$ (the concentration-cost relaxation
on the lower edge), with the second channel (criterion
devaluation, the C1-side $\CF$ effect) already pinned in
Section~\ref{sec:appendix-deriv-a1}.

\paragraph{Scope.}
Proposition~\ref{prop:r-dagger-rho} is a \emph{local}
statement about the lower edge of the escape band $(\rdagger(\val;\corr),
\rdagger(1;\corr))$; the peak location $r^\star(\val;\corr)$ is
not directly predicted. The argument assumes equicorrelated
Gaussian decision noise (the one-factor structure of
Section~\ref{sec:model-rho-channel}); structured non-equicorrelated
covariances are out of scope here. The numerical realisation runs
at the C2 headline cell under variant~A; variant~B is a mechanical
substitution of $\CR$ in \eqref{eq:c2-Kc-rho}--\eqref{eq:c2-Ku-rho}
and is left as a sensitivity probe
(Section~\ref{sec:limitations}). The conservation-family extension
to $p \ne 1$ inherits the $\corr = 0$ $\rdagger$-invariance of
Proposition~\ref{prop:r-dagger-invariance} of
Section~\ref{sec:extensions-a3} at the boundary
($\dprime_c = \dprime_u = \dprime_{\mathrm{base}}$ regardless of
$(\Rsens, p)$, so $K_c, K_u$ at $\corr = 0$ are $p$-independent);
the joint $(p, \corr > 0)$ band is left to a future increment.
The empirical drift in Table~\ref{tab:r-dagger-rho-drift} is
quantified at $\corr \in \{0, 0.2\}$, the central anchor of the
$\corr$ envelope of Section~\ref{sec:model-rho-channel}; a finer
$\corr$-grid bracketing of $d\rdagger/d\corr$ is queued as a
sensitivity refinement.

\paragraph{Reproducibility.}
The closed form \eqref{eq:r-dagger-rho} is implemented in
\texttt{Rebuild/derivations/verify\_C2\_rho/verify.py}. The
verification script (i) reconstructs $\rdagger(\val;0)$ on the
shared criterion grid and confirms recovery against the
$\corr = 0$ baseline of
Section~\ref{sec:results-c2} at
$\max|\Delta\rdagger| \le 5 \times 10^{-4}$ across
$\val \in \{1, 2, 3, 5, 8, 10\}$ (binary at $\val = 2$);
(ii) evaluates the drift table at $\corr = 0.2$
(Table~\ref{tab:r-dagger-rho-drift}) and confirms
sign-match against the empirical peak-drift signs of
Table~\ref{tab:rho-sensitivity} at $5/5$ rows
$\val \ne 1$; (iii) confirms the boundary-FD sign-flip
($\partial \Rew/\partial\alpha|_{1/\Nloc^{+}}$ switches sign across
$\Rsens = \rdagger(\val;\corr) \pm 0.05$) at $6/6$ probes
$(\val, \corr) \in \{1, 2, 3\} \times \{0, 0.2\}$. The numerical
payload is written to
\texttt{Rebuild/derivations/verify\_C2\_rho/output.json}; byte-identical
reruns produce the same SHA-256 digest
\texttt{ddbd3988e0253fde2cfae5906ca2749cda872d1382055b3d243b4b7c6a0678dc}.
The companion derivation file
\texttt{Rebuild/derivations/C2-{}-non-monotonic-vda-rho.md} (rb-023,
RB-026) is the long-form source of
\eqref{eq:c2-Ic}--\eqref{eq:r-dagger-rho} and includes the
additional algebraic steps elided here.

% ---------------------------------------------------------------------
\subsection{Symmetric recovery at $\Rsens = 1$}
\label{sec:appendix-c5}

The asymmetric benefit/cost family
$\benefit(\Rsens) = 2\Rsens/(\Rsens+1)$,
$\cost(\Rsens) = 2/(\Rsens+1)$ collapses at $\Rsens = 1$ to the
symmetric special case $\benefit = \cost = 1$ in which a single
shared transfer function drives both the cued and the uncued
sensitivity. The published manuscript states this collapse as a
numerical observation (``maximum difference: $0.0$ across all $210$
matched parameter combinations''). Here we record the claim at the
strength the rebuilt model actually licenses: a \emph{real-number
identity} that holds universally and a \emph{bit-exact} guarantee that
holds at the validation configuration. The treatment is a light-touch
consistency result; no new sweep is needed because the
$\corr \to 0$ recovery contract of
Section~\ref{sec:appendix-deriv-a1} already exercises $\Rsens = 1$ as
one of its three pinned cells (see Proposition~\ref{prop:c5-realnumber}
and the witness below).

\paragraph{The collapse as a real-number identity (universal claim).}
At $\Rsens = 1$ the two weights $\benefit(1) = 2/(1+1) = 1$ and
$\cost(1) = 2/(1+1) = 1$ are both \emph{exactly} representable in
IEEE-754 binary64 (the literal float $1.0$, hex \texttt{0x1.0p+0}).
The asymmetric sensitivity rule of Section~\ref{sec:model} writes
$\dprime_i(\alpha) =
 \dprime_{\mathrm{base}} +
 s_i(\Rsens)\,(\dprime_{\max} f(\alpha) - \dprime_{\mathrm{base}})$
where $s_i \in \{\benefit, \cost\}$ depending on whether location $i$
is attended or not; at $\Rsens = 1$ both slopes are unit, so the
expression collapses to $\dprime_i(\alpha) = \dprime_{\max} f(\alpha)$
--- the symmetric rule. Because the downstream signal-detection,
expected-reward, and criterion-optimisation stack is a deterministic
function of the $\dprime$ arrays alone, identical arrays force
identical optimal $(\alphacued^\star, c_c^\star, c_u^\star)$ and
identical $\Rew^\star$.

\begin{proposition}[Real-number recovery at $\Rsens = 1$]
\label{prop:c5-realnumber}
For every $(\valid, \val, \Nloc, \dprimemax, f_0, h)$ in the model's
domain, and at any cross-location correlation $\corr \in [0, 1)$, the
asymmetric model evaluated at $\Rsens = 1$ produces the same
$\dprime$ arrays as the symmetric special case
$\benefit = \cost = 1$. Consequently the optimal cued share, the
optimal cued and uncued criteria, and the optimal reward all match
the symmetric baseline as real-number identities:
$\alphacued^\star\!\big|_{\mathrm{asym},\,\Rsens=1} =
 \alphacued^\star\!\big|_{\mathrm{sym}}$,
$\Rew^\star\!\big|_{\mathrm{asym},\,\Rsens=1} =
 \Rew^\star\!\big|_{\mathrm{sym}}$, and
$\CF\!\big|_{\Rsens=1}$ matches the symmetric criterion-fraction
baseline pointwise. The cancellation
$\dprime_{\mathrm{base}} - \dprime_{\mathrm{base}} = 0$ is the whole
content of the word ``reduces.''
\end{proposition}

This is a property of the algebra, not of the numerical
implementation, and is therefore the appropriate \emph{universal}
form of the C5 claim.

\paragraph{Bit-exact $0.0$ at the validation configuration (a
Sterbenz-lemma band phenomenon).}
The literal ``maximum difference: $0.0$'' of the published
Appendix~A is sharper than Proposition~\ref{prop:c5-realnumber}: it
asserts not merely a real-number identity but a \emph{bit-for-bit}
identity under finite-precision arithmetic. We record it at its
proper scope. The asymmetric code path at $\Rsens = 1$ computes
$\mathrm{fl}\big(a + \mathrm{fl}(1.0 \cdot \mathrm{fl}(x - a))\big)$
with $a = \dprime_{\mathrm{base}}$ and
$x = \dprime_{\max} f(\alpha)$. Three observations compose:
(i) $\benefit(1)$ and $\cost(1)$ are the exact float $1.0$, so the
multiplicative slope round-trips identically;
(ii) Sterbenz's lemma~\cite{Sterbenz1974,Goldberg1991} states that
$a/2 \le x \le 2a$ implies $\mathrm{fl}(x - a) = x - a$ exactly;
(iii) the validation configuration of the published manuscript
($\Nloc = 4$, $\dprimemax = 2.0$, $f_0 = 0.5$, $h = \sqrt{\,\cdot\,}$,
variant~A) places
$a = \dprime_{\mathrm{base}} = \dprimemax(f_0 + (1{-}f_0) h(1/\Nloc))
 = 2.0\,(0.5 + 0.5 \cdot 0.5) = 1.5$,
and the swept output range $x = \dprimemax f(\alpha) \in [1.0, 2.0]$
sits inside the Sterbenz band $[a/2, 2a] = [0.75, 3.0]$ for every
$\alpha$ on the model's grid. The round-trip
$a + (x - a) = x$ is therefore exact at every grid point.

\begin{proposition}[Bit-exact symmetric recovery on the validation
configuration]
\label{prop:c5-sterbenz}
Under variant~A at the validation configuration above, the asymmetric
implementation evaluated at $\Rsens = 1$ returns
$\dprime$ arrays bit-identical to the symmetric implementation; hence
$\max\big|\Delta \alphacued^\star\big| = 0$ and
$\max\big|\Delta \Rew^\star\big| = 0$ exactly (not merely to machine
epsilon).
\end{proposition}

\paragraph{Scope: the literal $0.0$ is configuration-specific.}
Proposition~\ref{prop:c5-sterbenz}'s hypothesis is sufficient but not
necessary, and fails as soon as the smallest swept output
$\dprimemax f_0$ drops below $\dprime_{\mathrm{base}}/2$, which gives
the explicit threshold
\begin{equation}
\label{eq:c5-sterbenz-threshold}
f_0 \;<\; \frac{h(1/\Nloc)}{1 + h(1/\Nloc)}
\;=\; \frac{1}{3}
\qquad (h = \sqrt{\,\cdot\,},\; \Nloc = 4).
\end{equation}
Off-band configurations (notably $f_0 \lesssim 1/3$) lose bit-identity
and drift at the order of one unit in the last place ($\sim 10^{-17}$
to $10^{-16}$). The proper universal statement of C5 is therefore
\emph{``identical to machine precision''}; the literal ``$0.0$'' of
the published Appendix~A is a structural guarantee of the chosen
configuration, not a property of the model.

\paragraph{$\Rsens = 1$ is the smooth centre of the family, not a
knife-edge.}
The collapse at $\Rsens = 1$ is the limit of a smooth, two-sided
family rather than a removable singularity. The asymmetric slope ratio
$\benefit(\Rsens)/\cost(\Rsens) = \Rsens$ is continuous and
differentiable through $\Rsens = 1$, with
$d\benefit/d\Rsens|_{\Rsens=1} = 1/2$ and
$d\cost/d\Rsens|_{\Rsens=1} = -1/2$. The reviewer's continuity probe
at $\Rsens = 1 \pm \epsilon$ confirms the symmetric recovery is the
smooth interior of the family: $\max|\Delta \Rew^\star|$ against the
symmetric model scales linearly in $|\Rsens - 1|$, with slope
$\approx 0.084$ reward units per unit shift in $\Rsens$. Reporting
$\Rsens = 1$ as a special case is mathematically correct but
interpretively misleading if it suggests a knife-edge: it is the
\emph{balanced} point of the gain-modulation versus suppression
trade-off (see Section~\ref{sec:model}, Definition~\ref{def:three-levers}),
the natural null away from which the asymmetric VDA effects of
Section~\ref{sec:results} live.

\paragraph{C5 is conservation-form-invariant.}
The conservation rule on $(\benefit, \cost)$ is a model assumption,
not a derived statement (Section~\ref{sec:extensions-a3}). The
inherited additive constraint $\benefit + \cost = 2$ is the choice
$p = 1$ in the power-mean family $M_p(\benefit, \cost) = 1$ with
$\benefit/\cost = \Rsens$. At $\Rsens = 1$ the family identity
$\benefit/\cost = 1$ forces $\benefit = \cost$, and the constraint
$M_p(\benefit, \cost) = 1$ then forces $\benefit = \cost = 1$ for
every $p$ (every power mean of the constant sequence equals the
constant). Consequently the entire C5 argument
(Proposition~\ref{prop:c5-realnumber}--Proposition~\ref{prop:c5-sterbenz})
holds under any choice of conservation rule in the family without
modification: the recovery at $\Rsens = 1$ is
conservation-form-invariant by construction. The formal mechanism
is Proposition~\ref{prop:a3d-symmetric-corner} of
Section~\ref{sec:appendix-deriv-a3}; the C5-side corollary itself is
Corollary~\ref{cor:a3d-c5-invariance}. The
\emph{corollary} block of Section~\ref{sec:extensions-a3} records the
same statement from the A3 side; this paragraph closes the
cross-reference from the C5 side.

\paragraph{Reproducibility.}
The recovery contract is encoded in
\texttt{Rebuild/model/tests/test\_recovery.py} and is rerun under
\texttt{python3 -m model.tests.test\_recovery} from the
\texttt{Rebuild/} root. At the validation configuration and
$\corr = 0$, the policy-level test pins
$\VDA = 0.039825$, $\CF = 0.728228$, and
$\Rpone = 2.317239$ at $\Rsens = 1$ (variant~A, $\Nloc = 4$,
$\dprimemax = 2.0$, $f_0 = 0.5$, $h = \sqrt{\,\cdot\,}$), all
matched against the reviewer's logged numbers
(\texttt{Critique/replications/A1{-}{-}correlated{-}fa/output/results.json})
to a max absolute difference of $0.0$ on each of $\VDA$, $\CF$, and
$\Rpone$. The probability-level test confirms that the $\corr = 0$
quadrature shortcut returns the inherited independent product
$\Phinorm(b_c)\,\Phinorm(b_u)^{\Nloc - 1}$ of
Equation~\eqref{eq:pnofa-indep} to binary equality (max absolute
difference $0.0$). All seven recovery checks pass; the test driver
records a SHA-256 digest of
\texttt{d3c62215cedbe2f797c484a8d76e8bef1d5ac42cdb3092ff099a34aea1b6f18f}
on the numerical payload (output written to
\texttt{Rebuild/model/tests/recovery\_output.json}). The
conservation-form-invariance corollary above is independently
exercised by \texttt{Rebuild/model/tests/test\_conservation\_family.py}
(SHA-256 \texttt{f4f57a89e5108db47447cbeaf2f15440f56cdfffb65f3d173dc4a0550121791e}):
the two checks \texttt{test\_symmetric\_corner\_invariant} and
\texttt{test\_p0\_symmetric\_corner\_identity} verify
$\texttt{policies}(\Rsens{=}1, p{=}0) = \texttt{policies}(\Rsens{=}1, p{=}1)$
to floating-point identity across the variant grid, in agreement
with the closed-form $\benefit(1, p) = \cost(1, p) = 1$ identity.

% ---------------------------------------------------------------------
\subsection{Power-mean conservation family derivation}
\label{sec:appendix-deriv-a3}

This subsection promotes the conservation-family treatment of
Section~\ref{sec:extensions-a3} into a derivation in the rebuild's
voice. It (i) records the one-parameter power-mean family that
generalises the inherited additive A3 rule and re-derives its
closed-form weights; (ii) states the Hardy--Littlewood--Pólya
power-mean monotonicity inequality in exact pointwise form as a
KL divergence; (iii) proves the symmetric-corner identity that
backs the C5 conservation-form-invariance cross-reference of
Section~\ref{sec:appendix-c5}; (iv) supplies the full proof of
Proposition~\ref{prop:r-dagger-invariance} (the $\rdagger(\val)$
$p$-invariance) for which Section~\ref{sec:extensions-a3} gives
only a sketch; and (v) records the $d'$-channel chain rule for
$\partial \CF/\partial p$ that motivates the empirical
Theorem~\ref{thm:delta-cf-monotone} and isolates the one analytic
full-strength substatement that the chain rule yields. The
long-form source is the companion derivation file
\texttt{Rebuild/derivations/A3{-}{-}power-mean-conservation.md}
(rb-029, RB-033); this subsection records the equation skeleton,
the boxed closed forms, the two formal proofs, and the open-
conjecture status of $\Delta\CF \le 0$.

\paragraph{The power-mean conservation family.}
For $p \in \mathbb{R}$, the power (Hölder) mean of order $p$ on
positive reals $(x, y)$ is
\begin{equation}
  M_p(x, y) \;=\;
  \begin{cases}
    \displaystyle \left( \frac{x^p + y^p}{2} \right)^{1/p}, & p \ne 0, \\[4pt]
    \displaystyle \sqrt{x\,y}, & p = 0,
  \end{cases}
  \label{eq:a3d-power-mean}
\end{equation}
with limits $M_{-\infty}(x, y) = \min(x, y)$ and
$M_{+\infty}(x, y) = \max(x, y)$~\cite{HLP1934}. The A3
conservation family of Equation~\eqref{eq:conservation-family} is
the constraint $M_p(\benefit, \cost) = 1$ together with the
asymmetry ratio $\benefit/\cost = \Rsens$. Substituting
$\benefit = \Rsens\,\cost$ into $M_p(\benefit, \cost) = 1$ at
$p \ne 0$ gives
$1 = \cost\,\bigl((\Rsens^p + 1)/2\bigr)^{1/p}$ and hence the boxed
closed-form weights
\begin{equation}
  \boxed{\;
    \cost(\Rsens; p) \;=\;
    \left( \frac{2}{\Rsens^{\,p} + 1} \right)^{1/p},
    \qquad
    \benefit(\Rsens; p) \;=\;
    \Rsens\,\cost(\Rsens; p).
  \;}
  \label{eq:a3d-closed-form-weights}
\end{equation}
The geometric-mean limit $p \to 0$ reads
$\cost(\Rsens; 0) = 1/\sqrt{\Rsens}$,
$\benefit(\Rsens; 0) = \sqrt{\Rsens}$ (L'Hôpital on the inner log).
The cases $p = 1$ (additive, inherited paper) and $p = 0$
(multiplicative, the reviewer's attack form) recover the two scalar
rules: $\benefit(\Rsens; 1) = 2\Rsens/(\Rsens + 1)$,
$\cost(\Rsens; 1) = 2/(\Rsens + 1)$, and the geometric-mean limit
above. Equation~\eqref{eq:a3d-closed-form-weights} matches
Equation~\eqref{eq:beta-gamma-of-p} of
Section~\ref{sec:extensions-a3} and is implemented at
\texttt{Rebuild/model/core.py} as
\texttt{beta\_gamma(r, p)} (the $p = 1$ literal branch at line~311,
the $|p| < 10^{-12}$ geometric-mean branch at line~314, and the
general-$p$ branch at line~319). Three identities hold at every
$p \in \mathbb{R}$: (I-1) $\benefit/\cost = \Rsens$ by
construction; (I-2) $\benefit(1; p) = \cost(1; p) = 1$ (the
symmetric-corner identity, direct from
\eqref{eq:a3d-closed-form-weights} at $\Rsens = 1$); (I-3) the
sign of $\benefit - \cost$ is $\operatorname{sgn}(\Rsens - 1)$,
$p$-invariant.

\paragraph{HLP monotonicity in pointwise KL-divergence form.}
The classical Hardy--Littlewood--Pólya power-mean monotonicity
inequality~\cite{HLP1934} states that for fixed positive reals
$x \ne y$, the map $p \mapsto M_p(x, y)$ is strictly increasing in
$p$ over $\mathbb{R}$ (with $p = 0$ continuous by the
geometric-mean limit). Under the constrained form
$M_p(\benefit, \cost) = 1$ with $\benefit/\cost = \Rsens$, the
inequality translates into an exact pointwise identity on
$\partial \ln \cost/\partial p$. Let
$L(p) := \ln \cost(\Rsens; p) = (1/p)\,\ln\bigl(2/(\Rsens^p + 1)
\bigr)$ at $p \ne 0$. Direct differentiation and the substitution
of the Bernoulli weight
$\theta_p := \Rsens^p/(\Rsens^p + 1)$ give
\begin{equation}
  \boxed{\;
    \frac{\partial}{\partial p}\,\ln \cost(\Rsens; p)
    \;=\;
    -\frac{1}{p^2}\,
    \DKL\!\bigl(\,\mathrm{Bern}(\theta_p)\,\big\|\,
                    \mathrm{Bern}(1/2)\,\bigr),
    \qquad
    \theta_p \;=\; \frac{\Rsens^p}{\Rsens^p + 1}.
  \;}
  \label{eq:a3d-hlp-kl}
\end{equation}
The algebraic identity that the bracket of the direct expansion
$\ln(2/(\Rsens^p + 1)) + p\,\theta_p \ln \Rsens$ equals
$\theta_p \ln(2 \theta_p) + (1 - \theta_p)\,\ln(2(1 - \theta_p)) =
\DKL\!\bigl(\mathrm{Bern}(\theta_p) \,\|\, \mathrm{Bern}(1/2)\bigr)$
is given in the companion derivation file
\texttt{Rebuild/derivations/A3{-}{-}power-mean-conservation.md}
§3.2. A finite-difference cross-check of
\eqref{eq:a3d-hlp-kl} at seven test pairs
$(\Rsens, p) \in \{(0.3548, 0.5),\,(0.3548, 1.0),\,(10, 1.0),\,
(10, 0.5),\,(3.162, 2.0),\,(0.5, -1.0),\,(5.0, 1.5)\}$
(central FD, $\epsilon = 10^{-6}$) confirms LHS$-$RHS agreement
$\le 1.5 \times 10^{-10}$ in every case; at $\Rsens = 1$ both sides
are exactly $0$ (the Bernoulli reduces to $\mathrm{Bern}(1/2)$ and
$\DKL = 0$). Since $\DKL \ge 0$ with equality iff
$\theta_p = 1/2$ iff $\Rsens = 1$,
Equation~\eqref{eq:a3d-hlp-kl} yields, for every $\Rsens > 0$,
$\Rsens \ne 1$, every $p \in \mathbb{R} \setminus \{0\}$,
\begin{equation}
  \frac{\partial \cost}{\partial p}(\Rsens; p) \;<\; 0,
  \qquad
  \frac{\partial \benefit}{\partial p}(\Rsens; p) \;=\;
  \Rsens \cdot \frac{\partial \cost}{\partial p} \;<\; 0
  \quad (\Rsens > 0).
  \label{eq:a3d-mono-sign}
\end{equation}
Continuity through $p = 0$ holds by series expansion: with
$\theta_p = 1/2 + (p\,\ln \Rsens)/4 + O(p^2)$ and
$\DKL = 2(\theta_p - 1/2)^2 + O((\theta_p - 1/2)^3)$, the bracket
is $O(p^2)$ and the $1/p^2$ prefactor gives the finite limit
$\lim_{p \to 0}\partial \ln \cost/\partial p = -(\ln \Rsens)^2/8$.
The HLP envelope is therefore operative pointwise: as $p$
decreases, both weights $(\benefit, \cost)$ strictly increase at
every $\Rsens \ne 1$ — the conservation-form-sensitive side of A3,
witnessed empirically by the $+13.72\%$ shift of
$\cost(0.3548; 1) = 1.4762 \to \cost(0.3548; 0) = 1.6788$ at the
rb-016 $\val = 10$ peak $\Rsens$.

\paragraph{The symmetric corner and the C5 corollary.}
The companion derivation file factors out two formal statements
about the symmetric corner; we restate them here in the
manuscript's voice.

\begin{proposition}[Symmetric corner identity, $\Rsens = 1$ is
conservation-form-invariant]
\label{prop:a3d-symmetric-corner}
For every $p \in \mathbb{R}$ (including $p = 0$),
$\benefit(1; p) = \cost(1; p) = 1$.
\end{proposition}

\begin{proof}
At $\Rsens = 1$, $\Rsens^p = 1$ for every $p$, so
\eqref{eq:a3d-closed-form-weights} at $p \ne 0$ gives
$\cost = (2/2)^{1/p} = 1$ and $\benefit = 1 \cdot 1 = 1$. At
$p = 0$ the geometric-mean limit gives
$\cost = 1/\sqrt{1} = 1$ and $\benefit = 1$. \qed
\end{proof}

Proposition~\ref{prop:a3d-symmetric-corner} is an identity of the
conservation family alone — it does not depend on any other part
of the rebuilt model — and is the formal mechanism behind the C5
conservation-form-invariance cross-reference of
Section~\ref{sec:appendix-c5}.

\begin{corollary}[C5 symmetric recovery is conservation-form-invariant]
\label{cor:a3d-c5-invariance}
The C5 symmetric-recovery result of
Proposition~\ref{prop:c5-realnumber} (Section~\ref{sec:appendix-c5})
holds at every $p \in \mathbb{R}$. Specifically: at $\Rsens = 1$
and any $p$, the asymmetric $\dprime$-map collapses pointwise to
the symmetric baseline, $\dprime_c(\alpha; 1, p) =
\dprime_u((1 - \alpha)/(\Nloc - 1); 1, p) = \dprimemax\,f(\alpha)$
at every $\alpha$.
\end{corollary}

\begin{proof}
By Proposition~\ref{prop:a3d-symmetric-corner},
$\benefit(1; p) = \cost(1; p) = 1$ at every $p$. The rebuilt
$\dprime$-map at $\Rsens = 1$ then reads
$\dprime_c(\alpha; 1, p) = \dprime_{\mathrm{base}} + 1 \cdot
[\dprimemax\,f(\alpha) - \dprime_{\mathrm{base}}] = \dprimemax\,
f(\alpha)$, with no $p$-dependence anywhere on the right; the
uncued counterpart is identical with $(1 - \alpha)/(\Nloc - 1)$
in place of $\alpha$. The asymmetric model's per-policy expected
reward at $\Rsens = 1$ is therefore the symmetric baseline's,
giving the same $(\alpha^\star, c_c^\star, c_u^\star, \Rew^\star,
\CF)$ as real-number identities — Proposition~\ref{prop:c5-realnumber}.
\qed
\end{proof}

The numerical witnesses are the rb-015 family tests
\texttt{test\_symmetric\_corner\_invariant} and
\texttt{test\_p0\_symmetric\_corner\_identity} of
\texttt{Rebuild/model/tests/test\_conservation\_family.py} (SHA-256
\texttt{f4f57a89e5108db47447cbeaf2f15440f56cdfffb65f3d173dc4a0550121791e},
$14/14$ PASS), which verify
$\benefit(1, p) = \cost(1, p) = 1$ and
\texttt{policies(r=1, p=0) == policies(r=1, p=1)} to floating-point
identity across the variant grid. Combined with the rb-001
symmetric-recovery contract
\texttt{Rebuild/model/tests/test\_recovery.py} (SHA-256
\texttt{d3c62215cedbe2f797c484a8d76e8bef1d5ac42cdb3092ff099a34aea1b6f18f},
which pins $\VDA = 0.039825$ and $\CF = 0.728228$ at $\Rsens = 1$
to zero diff against the inherited $p = 1$ reference), this gives
bit-exact recovery at $\Rsens = 1$ across the conservation family
in the bit-exact band of Proposition~\ref{prop:c5-sterbenz}
(Sterbenz lemma).

\paragraph{Full proof of $\rdagger(\val)$ conservation-form invariance.}
Section~\ref{sec:extensions-a3} states the $p$-invariance of the
escape threshold $\rdagger(\val) = K_u(\val)/[(\Nloc - 1)\,
K_c(\val)]$ as Proposition~\ref{prop:r-dagger-invariance} with a
short proof sketch. We now record the full three-step proof, in
which every step ties to a specific structural property of the
$\dprime$-map at the uniform allocation.

\begin{proof}[Full proof of Proposition~\ref{prop:r-dagger-invariance}]
\textbf{(Step 1) Sensitivities at $\alpha = 1/\Nloc$ are
$p$-independent.} At the uniform allocation $\alpha = 1/\Nloc$,
the rebuilt $\dprime$-map reads
\begin{equation*}
  \dprime_c(1/\Nloc;\, \Rsens, p)
  \;=\;
  \dprime_{\mathrm{base}}
  + \benefit(\Rsens; p)
    \bigl[\,\dprimemax\,f(1/\Nloc)
            \,-\, \dprime_{\mathrm{base}}\,\bigr]
  \;=\;
  \dprime_{\mathrm{base}},
\end{equation*}
because the bracket $\dprimemax\,f(1/\Nloc) -
\dprime_{\mathrm{base}}$ vanishes by the definition
$\dprime_{\mathrm{base}} := \dprimemax\,f(1/\Nloc)$ of the
uniform-allocation baseline. The vanishing bracket annihilates
any value of $\benefit(\Rsens; p)$, so $\dprime_c$ at
$\alpha = 1/\Nloc$ is $\benefit$-independent — and hence
$p$-independent. At the same $\alpha = 1/\Nloc$ the per-uncued
allocation is $(1 - 1/\Nloc)/(\Nloc - 1) = 1/\Nloc$, so the
uncued bracket also vanishes and
$\dprime_u(1/\Nloc;\, \Rsens, p) = \dprime_{\mathrm{base}}$
independent of $\cost(\Rsens; p)$ and hence of $p$. The pair
$(\dprime_c, \dprime_u) = (\dprime_{\mathrm{base}},
\dprime_{\mathrm{base}})$ at $\alpha = 1/\Nloc$ is therefore
$p$-independent (and $\Rsens$-independent).

\textbf{(Step 2) $\textrm{P3}$-optimal criteria are $p$-independent.}
At the fixed uniform allocation $\alpha = 1/\Nloc$, the $\textrm{P3}$
expected-reward functional reduces to a function of
$(\dprime_c, \dprime_u, c_c, c_u; \valid, \val, \Nloc, \CR)$ whose
$p$-dependence enters only through $(\dprime_c, \dprime_u)$. By
Step 1, those are $p$-independent. The $\textrm{P3}$-optimal criteria
\begin{equation*}
  (c_c^\star, c_u^\star)
  \;=\;
  \operatorname*{arg\,max}_{(c_c, c_u)}\,
    \Rpthree\bigl(\dprime_{\mathrm{base}},
                       \dprime_{\mathrm{base}},
                       c_c, c_u;\,
                       \valid, \val, \Nloc, \CR\bigr)
\end{equation*}
are then the argmax of a functional whose every input is
$p$-independent; the criteria themselves are therefore
$p$-independent.

\textbf{(Step 3) $K_c, K_u$ are $p$-independent, hence $\rdagger$
is.} The partials
$K_c(\val) = \partial \Rpthree/\partial \dprime_c$ and
$K_u(\val) = \partial \Rpthree/\partial \dprime_u$ are evaluated
at the $p$-independent
$(\dprime_{\mathrm{base}}, \dprime_{\mathrm{base}}, c_c^\star,
c_u^\star)$, and the functional form of $\Rpthree$ depends only on
$(\valid, \val, \Nloc, \CR)$ — no conservation choice enters its
expression. Hence $K_c(\val)$ and $K_u(\val)$ are $p$-independent,
and so is their ratio
$\rdagger(\val) = K_u/[(\Nloc - 1)\,K_c]$. \qed
\end{proof}

The structural full proof is tighter than the empirical witness
allows: it says $\rdagger(\val)$ is an algebraic invariant of the
family, not merely a numerical coincidence. The rb-016
\texttt{recovery.test\_3\_r\_dagger\_p\_invariance} check
(payload digest
\texttt{055bf4ec862c711f2c6a3fa0831e1373b806a84c5f5c1a6b13fd1d9d7a039a33})
records $K_c$, $K_u$, $c_c^\star$, $c_u^\star$, and
$\rdagger(\val)$ identical to floating-point identity across
$p \in \{0,\, 1/2,\, 1\}$ and
$\val \in \{2,\, 3,\, 5,\, 8,\, 10\}$ (max $|\Delta| = 0.0$),
verifying that the algebra commutes with the floating-point
implementation across all three branches of \texttt{beta\_gamma}.
What is \emph{not} invariant is the peak inside the band: at
$\alpha \ne 1/\Nloc$ the bracket
$\dprimemax\,f(\alpha) - \dprime_{\mathrm{base}}$ no longer
vanishes, the $p$-dependence of $\benefit$ and $\cost$ does not
cancel, and the peak coordinates
$(\Rsens^\star, \VDA^\star)$ pick up the conservation-family band
of Table~\ref{tab:a3-c2-peak-band} and
Figure~\ref{fig:a3-vda-peak-band}. The C2 thread is therefore
characterised by the cleanest possible separation: the
\emph{left edge} of the positive-$\VDA$ band is a topological
invariant of the model, the \emph{peak} inside the band is a
calibration feature.

\paragraph{$d'$-channel chain rule for $\partial \CF/\partial p$.}
Theorem~\ref{thm:delta-cf-monotone} of
Section~\ref{sec:extensions-a3}
($\Delta\CF := \CF(p = 0) - \CF(p = 1) \le 0$ at every cell of
the $4{,}410$-cell C1 grid, with $0$ reverse flips) admits a
chain-rule sign analysis. With
$\CF = (\Rpthree - \Rpone)/(\Rpfour - \Rpone)$ from
\eqref{eq:cf-def-model} and the envelope theorem applied to each
policy reward (the optimum's first-order conditions remove the
inner gradient terms), the $p$-partials thread through the
$\dprime$-map at the policy's own
$(\alpha, c_c, c_u)$ optimum:
\begin{equation}
  \frac{\partial R_{\mathrm{P}_i}}{\partial p}
  \;=\;
  \frac{\partial R_{\mathrm{P}_i}}{\partial \dprime_c}\!\bigg|_{\!\star}
  \cdot
  \frac{\partial \dprime_c}{\partial p}
  \;+\;
  \frac{\partial R_{\mathrm{P}_i}}{\partial \dprime_u}\!\bigg|_{\!\star}
  \cdot
  \frac{\partial \dprime_u}{\partial p},
  \qquad
  i \in \{1, 3, 4\},
  \label{eq:a3d-chain-rule}
\end{equation}
where from the $\dprime$-map definition,
\begin{equation}
  \frac{\partial \dprime_c}{\partial p}
  \;=\;
  \frac{\partial \benefit}{\partial p}
  \,\bigl[\,\dprimemax\,f(\alpha) - \dprime_{\mathrm{base}}\,\bigr],
  \qquad
  \frac{\partial \dprime_u}{\partial p}
  \;=\;
  \frac{\partial \cost}{\partial p}
  \,\bigl[\,\dprimemax\,f((1-\alpha)/(\Nloc-1))
           - \dprime_{\mathrm{base}}\,\bigr].
  \label{eq:a3d-dprime-of-p}
\end{equation}
By Equation~\eqref{eq:a3d-mono-sign}, both
$\partial \benefit/\partial p$ and $\partial \cost/\partial p$ are
strictly negative at every $\Rsens \ne 1$, every $p$. For
$\alpha \ge 1/\Nloc$ (the cued-enhancement regime of the
$\textrm{P}$ family) the first bracket of
\eqref{eq:a3d-dprime-of-p} is non-negative and the second is
non-positive, so $\partial \dprime_c/\partial p \le 0$ and
$\partial \dprime_u/\partial p \ge 0$ — the cued and uncued $d'$
channels carry opposite sign of motion in $p$. Substituting into
\eqref{eq:a3d-chain-rule}, the two channel contributions to
$\partial R_{\mathrm{P}_i}/\partial p$ are of opposite sign; the resolution per
policy depends on the relative weights
$(\partial R_{\mathrm{P}_i}/\partial \dprime_c,\, \partial R_{\mathrm{P}_i}/\partial
\dprime_u)$ at the policy's own optimum.

\begin{proposition}[$\Rpthree$ at $\alpha = 1/\Nloc$ is $p$-invariant]
\label{prop:a3d-rp3-invariance}
At the $\textrm{P3}$ allocation $\alpha = 1/\Nloc$,
$\partial \Rpthree/\partial p \equiv 0$.
\end{proposition}

\begin{proof}
$\textrm{P3}$ is the criterion-only policy fixed at $\alpha = 1/\Nloc$.
By Step 1 of the proof of
Proposition~\ref{prop:r-dagger-invariance} above, the brackets in
\eqref{eq:a3d-dprime-of-p} vanish at $\alpha = 1/\Nloc$, so
$\partial \dprime_c/\partial p = \partial \dprime_u/\partial p =
0$, and \eqref{eq:a3d-chain-rule} gives
$\partial \Rpthree/\partial p = 0$. \qed
\end{proof}

Proposition~\ref{prop:a3d-rp3-invariance} is the one analytic
full-strength substatement that the chain rule yields. $\textrm{P1}$
optimises $\alpha^\star$ away from $1/\Nloc$ and $\textrm{P4}$
optimises both $\alpha$ and the criteria; neither enjoys the
vanishing-bracket collapse, and both pick up the opposite-sign
$d'$-channel competition above. The combined expression for the
$\CF$ gradient,
\begin{equation*}
  \frac{\partial \CF}{\partial p}
  \;=\;
  \frac{(\Rpfour - \Rpone)\,(\partial_p \Rpthree
                              - \partial_p \Rpone)
       \;-\; (\Rpthree - \Rpone)\,(\partial_p \Rpfour
                              - \partial_p \Rpone)}
       {(\Rpfour - \Rpone)^2},
\end{equation*}
collapses by Proposition~\ref{prop:a3d-rp3-invariance} to
\begin{equation*}
  \frac{\partial \CF}{\partial p}
  \;=\;
  \frac{-(\Rpfour - \Rpone)\,\partial_p \Rpone
       \;-\; (\Rpthree - \Rpone)\,(\partial_p \Rpfour
                              - \partial_p \Rpone)}
       {(\Rpfour - \Rpone)^2}.
\end{equation*}
Both $\Rpfour - \Rpone$ and $\Rpthree - \Rpone$ are positive
(definition of $\CF$); the sign of $\partial \CF/\partial p$
therefore reduces to a competition between
$-\partial_p \Rpone$ (first term, sign $\ge 0$ since
$\partial_p \Rpone \le 0$) and
$-(\partial_p \Rpfour - \partial_p \Rpone)$ (second term). The
uniform inequality
$|\partial_p \Rpfour| \le |\partial_p \Rpone|$ (the joint policy
falls in $p$ no faster than the sensitivity-only policy) suffices
to make $\partial \CF/\partial p \le 0$ pointwise, which would
integrate to $\Delta\CF \le 0$. The inequality holds at $0$
reverse flips out of $4{,}410$ cells in the rb-016 Block B sweep,
but a closed-form proof has so far eluded us: both
$\partial_p \Rpone$ and $\partial_p \Rpfour$ pick up
policy-specific $\alpha^\star(p)$ moving-target dependencies that
the chain rule does not collapse. The rebuild therefore reports
Theorem~\ref{thm:delta-cf-monotone} as \emph{empirical at full
strength} (the $4{,}410$-cell statement, $0$ reverse flips), with
the chain-rule analysis \eqref{eq:a3d-chain-rule}--%
\eqref{eq:a3d-dprime-of-p} identifying the sign mechanism and
Proposition~\ref{prop:a3d-rp3-invariance} contributing the one
analytic invariant, and leaves the uniform closed-form proof as a
queued open problem (see the \emph{Scope} paragraph below and the
companion derivation file §8.2).

\paragraph{Scope.}
The derivation is held at $\corr = 0$ throughout: the joint
$(p, \corr)$ sensitivity would compose the HLP gradient
\eqref{eq:a3d-hlp-kl} with the $\corr$-channel Slepian gradient
of Section~\ref{sec:appendix-deriv-c2}, but is not in the
rb-016 Block B sweep that supplies the numerical witnesses
(\texttt{Rebuild/sims/A3{-}{-}conservation-band/output} digest
\texttt{055bf4ec862c711f2c6a3fa0831e1373b806a84c5f5c1a6b13fd1d9d7a039a33}).
The derivation is also held at the inherited per-location
homogeneous $\Rsens$ (A2) and homogeneous uncued allocation (A8);
both are relaxed in Sections~\ref{sec:extensions-a2}
and~\ref{sec:extensions-a8} respectively, and the relaxations
compose linearly with $p$ (the conservation-family weights
$(\benefit, \cost)$ enter the per-location $\dprime$-map
independently of the heterogeneity structure). A closed-form
algebraic proof of the per-cell $\Delta\CF \le 0$ monotonicity on
the $4{,}410$-cell grid would require simultaneous control of
$\partial_p \Rpone$ and $\partial_p \Rpfour$ along the
$\alpha^\star(p)$ moving target, and is not yet in hand.

\paragraph{Reproducibility.}
Equation~\eqref{eq:a3d-closed-form-weights} is exercised by
\texttt{Rebuild/model/tests/test\_conservation\_family.py}
(SHA-256
\texttt{f4f57a89e5108db47447cbeaf2f15440f56cdfffb65f3d173dc4a0550121791e};
$14/14$ PASS). The
\texttt{test\_family\_identities} check verifies $\benefit/\cost
= \Rsens$ and $M_p(\benefit, \cost) = 1$ to relative error
$\le 4.4 \times 10^{-16}$ across
$p \in \{-2, -1, -1/2, 0, 1/2, 1, 2\}$ and a $21$-point
log-$\Rsens$ grid $\Rsens \in [0.1, 10]$, backing
\eqref{eq:a3d-closed-form-weights} pointwise.
The
\texttt{test\_symmetric\_corner\_invariant} and
\texttt{test\_p0\_symmetric\_corner\_identity} checks back
Proposition~\ref{prop:a3d-symmetric-corner} and
Corollary~\ref{cor:a3d-c5-invariance} to floating-point identity.
Proposition~\ref{prop:r-dagger-invariance} is exercised by the
rb-016 \texttt{recovery.test\_3\_r\_dagger\_p\_invariance} check
(payload digest
\texttt{055bf4ec862c711f2c6a3fa0831e1373b806a84c5f5c1a6b13fd1d9d7a039a33};
max $|\Delta| = 0.0$). The finite-difference cross-check of
\eqref{eq:a3d-hlp-kl} at the seven test pairs above (LHS$-$RHS
$\le 1.5 \times 10^{-10}$ in every case) reruns at the companion
derivation file
\texttt{Rebuild/derivations/A3{-}{-}power-mean-conservation.md} §3.2
and is the long-form source of the equation labels, the boxed
identities, and the two formal proofs in this subsection.
